\section{Four instruments}\label{sec:design}

Each arm is specified precisely enough to cost and, in the
microsimulation stage, to legislate hypothetically inside the tax--benefit
model. All four share two eligibility definitions. \emph{Narrow}:
employees made redundant from manufacturing (SIC section C)---the
observable proxy for trade exposure, since certification of trade
causation, as in the US petition system, has no UK analogue.
\emph{Broad}: all employees made redundant, any industry---the
administrable rule that avoids certification entirely. Any real scheme
sits between these bounds; a memo row prices a trade-exposed subset of
manufacturing at 30 per cent of the flow---a placeholder share, since
the companion study's public artifacts do not report the exposed-division
employment share.

\paragraph{Arm A: wage insurance.} A worker re-employed within a year of
redundancy at lower pay receives 50 per cent of
the gap between prior and new earnings, paid alongside wages for up to
two years, capped at prior earnings. This is the Reemployment TAA design
evaluated by \citet{hyman2024}, transplanted with two UK adaptations: no
age floor (their age-50 rule was a budget device, not a design
principle), and payments treated as taxable earnings so the instrument
interacts with income tax, National Insurance and the Universal Credit
taper---an interaction absent from the US programme's setting and
quantifiable only in the microsimulation stage. The arm is not foreign to
UK statute: ISERBS paid redundant steelworkers a weekly make-up to 90 per
cent of previous earnings on re-employment \citep{iserbs1988}, a more
generous replacement rate than costed here; Arm A is, in effect, that
scheme generalised beyond steel and detached from its expired European
funding.

\paragraph{Arm B: Unemployment Insurance for displaced workers.} The
government's proposed Unemployment Insurance---a single time-limited
contributory benefit at \pounds\TaaUiWeekly\ a week (2026--27 basis) replacing New
Style JSA and ESA---applied to redundant workers. The costed quantity is
\emph{incremental}: the difference between the proposed benefit and the
current New Style JSA baseline (\pounds\TaaJsaWeekly\ a week for 26
weeks), at the government's candidate durations, with every claimant assumed to
draw the full entitlement---an exhaustion convention that makes the
costed figure an upper bound, since claimants re-employed inside the
benefit window draw only part of it. This arm is deliberately
parasitic on a live reform: it asks what the reform the government has
already proposed would deliver specifically to displaced workers, and
what topping its duration up to the US-style horizon would add.

\paragraph{Arm C: retraining stipend.} A flat
\pounds\TaaTrainPerWorker\ payment to redundant workers entering
approved retraining. The US evidence expects little from this arm
\citep{damicoschochet2012}; it is included because every real-world
adjustment package contains it, and because its cost--coverage profile
(cheap per head, take-up-limited) contrasts instructively with the
insurance arms.

\paragraph{Arm D: redundancy-pay disregard in Universal Credit.}
Statutory redundancy pay counts toward the Universal Credit capital test,
whose thresholds (\pounds 6{,}000--\pounds 16{,}000) have been frozen in
cash terms since 2006: a redundancy payment can extinguish the very
means-tested support displacement triggers. This arm disregards
redundancy pay from the capital test for twelve months after job loss.
Its gross cost cannot be computed from public data---it depends on the
joint distribution of redundancy pay, household capital and UC
entitlement, observable only in microdata---so Section~\ref{sec:results}
reports a labelled exposure bound and the microsimulation stage carries
the real costing. It is nonetheless the arm most directly implied by the
companion study, which found the UC capital gate among the unmodelled
channels most likely to reduce displacement-side cushioning in practice.

Two arms pay people to work (A), or to have worked (B); one pays them to
retrain (C); one stops a lump sum from switching off their safety net
(D). The set spans the design space the US literature argues over, and
each arm maps onto a distinct behavioural risk that the costing brackets
rather than resolves: take-up (all arms), re-employment response (A),
duration response (B), training deadweight (C), and employer gaming of
redundancy structuring (D).
